\documentclass[11pt,aspectratio=169]{beamer}
\usepackage[T1]{fontenc}
\usepackage{lmodern}
\usepackage{amsmath,amssymb,booktabs,array}
\usetheme{default}
\definecolor{navy}{RGB}{20,50,110}
\definecolor{brick}{RGB}{160,40,30}
\setbeamercolor{structure}{fg=navy}
\setbeamercolor{title}{fg=navy}
\setbeamercolor{frametitle}{fg=navy}
\setbeamercolor{alerted text}{fg=brick}
\setbeamercolor{block title}{fg=white,bg=navy}
\setbeamercolor{block body}{bg=navy!5}
\setbeamerfont{frametitle}{series=\bfseries,size=\large}
\setbeamerfont{title}{series=\bfseries,size=\Large}
\setbeamertemplate{navigation symbols}{}
\setbeamertemplate{footline}{\hfill\scriptsize\textcolor{gray}{
  \insertframenumber\,/\,\inserttotalframenumber}\hspace{3mm}\vspace{2mm}}
\setbeamertemplate{itemize items}[circle]
\setbeamertemplate{blocks}[default]
\setbeamersize{text margin left=9mm,text margin right=9mm}
\setlength{\parskip}{5pt}
\newcommand{\N}{\mathbb{N}}
\newcommand{\Z}{\mathbb{Z}}
\newcommand{\Q}{\mathbb{Q}}
\newcommand{\R}{\mathbb{R}}
\newcommand{\stir}[2]{\genfrac{\{}{\}}{0pt}{}{#1}{#2}}
\newcommand{\ind}[1]{\mathbf{1}_{[#1]}}
\newcommand{\lean}[1]{\texttt{\detokenize{#1}}}
\newcommand{\takeaway}[1]{\par\medskip{\color{navy}\bfseries #1}\par}
\newcommand{\source}[2]{\par\medskip{\scriptsize Source: \href{#1}{#2}.}}
\title{Erd\H{o}s Problem 252}
\subtitle{Factorial tails, exact cancellation, and one surviving shift}
\author{Tokengrinder}
\date{September 2026}
\hypersetup{pdftitle={Erdos 252: explaining the proof},pdfauthor={Tokengrinder},
  pdfsubject={A guided exposition of the Lean proof},
  pdfkeywords={Erdos 252, irrationality, factorial series, Lean}}
\begin{document}
\begin{frame}[plain]
  \titlepage
\end{frame}
\begin{frame}{The theorem concerns every divisor-power sum}

For $n\ge1$ and $k\in\N=\{0,1,2,\ldots\}$, define
\[
 \sigma_k(n)=\sum_{d\mid n}d^k,
 \qquad
 \alpha_k=\sum_{n\ge1}\frac{\sigma_k(n)}{n!}.
\]
For example, the divisors of $6$ are $1,2,3,6$:
\[
 \sigma_0(6)=4,\qquad \sigma_1(6)=12,\qquad \sigma_2(6)=50.
\]
(The series converges: $\sigma_k(n)\le n^{k+1}$, and factorial growth wins.)
\begin{block}{Main theorem}
$\alpha_k$ is irrational for every natural number $k$.
\end{block}
We explain the positive degrees first; $k=0$ is exercise.

\note{[Sources] Local proof: Erdos252/Solution.lean and PROOF.tex. [/Sources]}
\end{frame}

% ----------------------------------------------------------------------

\begin{frame}{Earlier results reached the fourth power}

\begin{center}
\begin{tabular}{@{}lp{9.7cm}@{}}
\toprule
Exponent & Prior results \\
\midrule
$k=0$ & Elementary factorial-tail argument.\\[2pt]
$k=1,2$ & Classical cases; see Pratt's 1950s references [9,11].\\[2pt]
$k=3$ & Schlage-Puchta (2006) and Friedlander--Luca--Stoiciu (2007),
independently.\\[2pt]
$k=4$ & Pratt (2022).\\[2pt]
$k\ge5$ & now solved\\
\bottomrule
\end{tabular}
\end{center}
Earlier all-degree results assumed Schinzel's hypothesis H or a suitable
prime-tuples conjecture.
\takeaway{The argument here uses fixed coprime integers and exact cancellation.}
\source{https://arxiv.org/abs/2209.11124}{Pratt (2022), introduction and references}

\note{[Sources] https://arxiv.org/pdf/2209.11124, introduction;
https://math.colgate.edu/~integers/h31/h31.pdf, introduction;
https://www.erdosproblems.com/forum/thread/252?order=newest,
catalogue revision dated 22 January 2026.
[/Sources] This is a description of prior literature, not a claim that an external
catalogue has accepted this formal artifact. Pratt's references [9,11] distinguish
the 1952/1953 problem statements from the 1953/1954 published solutions.}
\end{frame}

% ----------------------------------------------------------------------

\begin{frame}{Warm-up: factorials expose a rational denominator}

Suppose $e=\sum_{m\ge0}1/m!=a/b$ with $a,b\ge1$ integers.
Then
\[
 b!\,e-b!\sum_{m\le b}\frac1{m!}
   =b!\sum_{m>b}\frac1{m!}\ \in\Z .
\]
Yet
\[
 0<b!\sum_{m>b}\frac1{m!}
   <\sum_{j\ge1}\frac1{(b+1)^j}
   =\frac1b\le1.
\]
There is no integer strictly between $0$ and $1$.
\begin{block}{The principle we will reuse}
A sequence of integers tending to $0$ is eventually exactly $0$.
\end{block}

\note{[Sources] Local proof: Erdos252/Solution.lean and PROOF.tex. [/Sources]}
\end{frame}

% ----------------------------------------------------------------------

\begin{frame}{Rationality would make all late tails integers}

Fix $k$. Scale the tail after the $n$-th term:
\[
 T_n:=n!\sum_{m>n}\frac{\sigma_k(m)}{m!}.
\]
If $\alpha_k=a/b$ with $a\in\Z$ and $b\ge1$, then for $n\ge b$,
\[
 T_n=
 \underbrace{\frac{n!a}{b}}_{\text{integer}}
 -
 \underbrace{\sum_{m=1}^{n}\sigma_k(m)\frac{n!}{m!}}_{\text{integer}}.
\]
All later work will rely on:
\[
 \alert{\alpha_k\in\Q\quad\Longrightarrow\quad T_n\in\Z
 \text{ for all sufficiently large }n.}
\]

\note{[Sources] Local proof: eventually_scaledTail_integral and
seriesPrefix_integral. [/Sources] The talk's T_n equals Lean's scaledTail k (n+1);
this indexing convention is used consistently throughout.}
\end{frame}

% ----------------------------------------------------------------------

\begin{frame}{The original tails need not become small}

The first terms are
\[
 T_n=\frac{\sigma_k(n+1)}{n+1}
 +\frac{\sigma_k(n+2)}{(n+1)(n+2)}+\cdots .
\]
The divisor $n+1$ alone gives the lower bound
\[
 T_n\ge\frac{\sigma_k(n+1)}{n+1}\ge(n+1)^{k-1}\qquad(k\ge1).
\]
\begin{itemize}
\item For $k=1$, the tail is at least $1$; it cannot tend to $0$.
\item For $k\ge2$, this lower bound tends to infinity.
\end{itemize}
\takeaway{We need a small integer combination of tails, not a small individual tail.}

\note{[Sources] Local proof: Erdos252/Solution.lean and PROOF.tex. [/Sources]}
\end{frame}

% ----------------------------------------------------------------------

\begin{frame}{The proof has four main moves}

\begin{enumerate}
\item \textbf{Expand.} Separate a finite main term from an error smaller than $1/n$.
\medskip
\item \textbf{Align.} Read several tails at compatible dilated indices.
\medskip
\item \textbf{Cancel.} Finite differences remove all lower powers.
The weighted integer is eventually zero, forcing the remaining sum to tend to zero.
\medskip
\item \textbf{Separate.} Two progression averages detect one surviving shift
differently. They cannot both be zero.
\end{enumerate}

\note{[Sources] Local proof: Erdos252/Solution.lean and PROOF.tex. [/Sources]}
\end{frame}

% ----------------------------------------------------------------------

\begin{frame}{1. Reciprocal products have integer coefficients}

For all suitably large $x$ we have:
\[
\begin{aligned}
 \frac1{x(x-1)}&=\frac1{x^2}+\frac1{x^3}+\frac1{x^4}+\cdots,\\
 \frac1{x(x-1)(x-2)}&=\frac1{x^3}+\frac3{x^4}+\frac7{x^5}+\cdots.
\end{aligned}
\]
In general, for $x>j$,
\[
 \frac1{x(x-1)\cdots(x-j)}
   =\sum_{i\ge j}\stir{i}{j}\frac1{x^{i+1}}.
\]
$\stir{i}{j}$ counts partitions of $i$ objects into $j$ nonempty blocks.
In particular, $\stir{i}{j}=0$ for $i<j$, and $\stir{j}{j}=1$.
\takeaway{We will only keep the powers through $x^{-(k+1)}$.}

\note{[Sources] Local proof: stirlingPoly_recurrence, stirlingErr_recurrence,
stirlingErr_bounds. [/Sources] The infinite generating function is an explanatory identity.
Lean proves the finite expansion and its error directly by recurrence; no unproved
infinite-series expansion is imported.}
\end{frame}

% ----------------------------------------------------------------------

\begin{frame}{1. The discarded part is smaller than $1/n$}

With $j$ indexing the shifted argument and $i+1$ the denominator power,
\[
 T_n=\sum_{j=0}^{k}\sum_{i=j}^{k}\stir{i}{j}
       \frac{\sigma_k(n+j+1)}{(n+j+1)^{i+1}}+E_n.
\]
The error has two parts:
\begin{itemize}
\item the terms after the first $k+1$ terms of the actual tail;
\item the higher reciprocal powers of each retained denominator.
\end{itemize}
Divisor pairing gives $\sigma_k(m)\le64m^k\sqrt m$.
A geometric majorant and the finite Stirling remainder bound then give
\[
 0\le E_n\le C_k n^{-3/2}\quad\text{for large }n,
 \qquad\alert{nE_n\longrightarrow0}.
\]
Here $C_k$ is a fixed constant depending only on $k$.

\note{[Sources] Local proof: Erdos252/Solution.lean and PROOF.tex. [/Sources]}
\end{frame}

% ----------------------------------------------------------------------

\begin{frame}{2. Multiplicativity turns dilations into shifts}

Choose a positive integer $p$ coprime to $1,\ldots,k+1$, and an offset $0\le t<p$.
For large $N\equiv t\pmod{p^2}$, put $n=(N-t)/p$.
\[
 p\mid n,\qquad
 \gcd(p,n+j+1)=\gcd(p,j+1)=1 \quad(0\le j\le k).
\]
Define $r=(j+1)p-t$. Then $p(n+j+1)=N+r$, so
\[
 \sigma_k(p)\,\sigma_k(n+j+1)=\sigma_k(N+r).
\]
Consequently,
\[
 \boxed{\ \sigma_k(p)\frac{\sigma_k(n+j+1)}{(n+j+1)^{i+1}}
       =p^{i+1}\frac{\sigma_k(N+r)}{(N+r)^{i+1}}\ }.
\]

\note{[Sources] Local proof: Erdos252/Solution.lean and PROOF.tex. [/Sources]}
\end{frame}

% ----------------------------------------------------------------------

\begin{frame}{2. Example: three tails share the same first shift}

For $k=1$, use three pairwise coprime multipliers, including the composite $9$:
\begin{center}
\renewcommand{\arraystretch}{1.35}
\begin{tabular}{@{}c|ccc@{}}
 & first tail & second tail & third tail\\
\hline
$p$ & $5$ & $7$ & $9$\\
$t$ & $0$ & $2$ & $4$\\
$n_p=(N-t)/p$ & $N/5$ & $(N-2)/7$ & $(N-4)/9$\\
$p(n_p+1)$ & $N+5$ & $N+5$ & $N+5$\\
$p(n_p+2)$ & $N+10$ & $N+12$ & $N+14$
\end{tabular}
\end{center}
Choose $N\equiv0\ (25)$, $N\equiv2\ (49)$, $N\equiv4\ (81)$.
One explicit solution progression is $N=47875+99225u$, $u\in\N$.

\note{[Sources] Local proof: Erdos252/Solution.lean and PROOF.tex. [/Sources]}
\end{frame}

% ----------------------------------------------------------------------

\begin{frame}{3. Finite differences cancel polynomial factors}

For every polynomial $P$ of degree at most $k$,
\[
 \sum_{a=0}^{k+1}b(a)P(a)=0,
 \qquad b(a)=(-1)^{k+1-a}\binom{k+1}{a}.
\]
For $k=1$, this is the second-difference identity:
\[
 P(0)-2P(1)+P(2)=0.
\]
In our example $P(a)=5+2a$, so
\[
 \left(\underbrace{5-2\cdot7+9}_{0}\right)\,
       \frac{\sigma_1(N+5)}{N+5}=0.
\]

\note{[Sources] Local proof: Erdos252/Solution.lean and PROOF.tex. [/Sources]}
\end{frame}

% ----------------------------------------------------------------------

\begin{frame}{3. Build one free coordinate for each lower shift}

Let $e=(e_0,\ldots,e_{k-1})\in\{0,\ldots,k+1\}^k$.
Set
\[
 G=(k+2)^k-1,\qquad B=1+k\,G!\,G,
\]
and define
\[
 p_e=B+G!\sum_{a<k}e_a(k+2)^a,\qquad
 t_e=G!\sum_{a<k}(a+1)e_a(k+2)^a.
\]
Their \emph{shift} at position $j$ is
\[
 r_e(j)=(j+1)p_e-t_e
       =(j+1)B+G!\sum_{a<k}(j-a)e_a(k+2)^a.
\]
\takeaway{For $j<k$, the shift's coefficient for $e_j$ is zero.}

\note{[Sources] Local proof: Erdos252/Solution.lean and PROOF.tex. [/Sources]}
\end{frame}

% ----------------------------------------------------------------------

\begin{frame}{3. The grid also guarantees the needed coprimality}

Write $I_e=\sum_{a<k}e_a(k+2)^a$. Distinct vertices $e$ give distinct
integers $I_e\in[0,G]$, and $p_e=B+G!\,I_e\equiv1\pmod{G!}$.
\begin{itemize}
\item \textbf{Pairwise coprime.}
A common prime divisor $\ell$ of $p_e,p_f$ would divide
$|I_e-I_f|\in[1,G]$, but not $G!$.
Then $\ell\le G$ forces $\ell\mid G!$, a contradiction.
\item \textbf{Coprime to small shifts.}
Since $k+1\le G$, we have $1,\ldots,k+1$ divides $G!$ and is coprime to $p_e$.
\item \textbf{Positive arguments.}
$t_e\le k\,G!\,I_e\le k\,G!\,G=B-1<p_e$, so every $r_e(j)>0$.
\end{itemize}

\note{[Sources] Local proof: Erdos252/Solution.lean and PROOF.tex. [/Sources]}
\end{frame}

% ----------------------------------------------------------------------

\begin{frame}{3. A common progression preserves integrality}

Use the product weights $w_e=\prod_{a<k}b(e_a)\in\Z$.
Since the $p_e^2$ are pairwise coprime, choose
\[
 A\equiv t_e\pmod{p_e^2}\quad\text{for every }e,\qquad
 Q=\prod_e p_e^2.
\]
Every $N=A+Qu$ satisfies the same congruences.
For large $u$, let $n_e=(N-t_e)/p_e$ and form
\[
 \boxed{\ W(N)=\sum_e w_e\,\sigma_k(p_e)\,T_{n_e}\ }.
\]
If $\alpha_k$ is rational, every $T_{n_e}$ is eventually integral:
\[
 \alert{W(A+Qu)\in\Z\quad\text{for all sufficiently large }u.}
\]

\note{[Sources] Local proof: Erdos252/Solution.lean and PROOF.tex. [/Sources]}
\end{frame}

% ----------------------------------------------------------------------

\begin{frame}{3. Only the top denominator power remains}

After expansion and using part 2 we have
\[
 W(N)=\sum_{j=0}^k\sum_{i=j}^k\stir{i}{j}\sum_e
      w_e p_e^{i+1}\frac{\sigma_k(N+r_e(j))}{(N+r_e(j))^{i+1}}+E(N).
\]
For every $i<k$:
\begin{itemize}
\item If $j<k$, hold all coordinates except $e_j$ fixed.
The shifted value is constant in $e_j$; the polynomial degree is
$i+1\le k$. Its weighted sum is zero.
\item If $j=k$, the Stirling coefficient $\stir{i}{k}$ is zero.
\end{itemize}
\takeaway{All rows $i<k$ vanish identically. Only $i=k$ survives.}

\note{[Sources] Local proof: Erdos252/Solution.lean and PROOF.tex. [/Sources]}
\end{frame}

% ----------------------------------------------------------------------

\begin{frame}{3. Name what survives before taking any limits}

Define the normalized divisor sum and the remaining coefficients:
\[
 \rho_k(m)=\frac{\sigma_k(m)}{m^k},\qquad
 c_{e,j}=w_e p_e^{k+1}\stir{k}{j}.
\]
Then cancellation gives
\[
 W(N)=S(N)+E(N),\qquad
 S(N):=\sum_e\sum_{j=0}^k
      c_{e,j}\frac{\rho_k(N+r_e(j))}{N+r_e(j)}.
\]
Here $E(N)=\sum_e w_e\sigma_k(p_e)E_{n_e}$ is a signed error.
For each fixed vertex $e$
\[
 NE_{n_e}=(n_e p_e+t_e)E_{n_e}\longrightarrow0.
\]
The family is fixed and finite; therefore $\alert{NE(N)\to0}$ as $u\to\infty$.

\note{[Sources] Local proof: Erdos252/Solution.lean and PROOF.tex. [/Sources]}
\end{frame}

% ----------------------------------------------------------------------

\begin{frame}{First limit: the weighted integer is eventually zero}

All limits here are along the same progression $N=A+Qu$.
From the divisor bound,
\[
 0\le\frac{\rho_k(m)}m\le\frac{64}{\sqrt m}\longrightarrow0.
\]
There are only finitely many fixed coefficients and shifts, so
\[
 S(N)\longrightarrow0,\qquad E(N)\longrightarrow0.
\]
Thus $W(N)=S(N)+E(N)\to0$.
\begin{block}{Use integrality at exactly this point}
Eventually $W(N)\in\Z$ and $|W(N)|<1$.
Hence $W(N)=0$, and consequently $S(N)=-E(N)$.
\end{block}

\note{[Sources] Local proof: Erdos252/Solution.lean and PROOF.tex. [/Sources]}
\end{frame}

% ----------------------------------------------------------------------

\begin{frame}{Second limit: remove the remaining denominators}

Define the \alert{survivor}, with exactly the same coefficients and shifts:
\[
 R(N):=\sum_e\sum_{j=0}^k c_{e,j}\rho_k(N+r_e(j)).
\]
For each fixed term, write $r=r_e(j)$ and $c=c_{e,j}$. Then
\[
 N\,\frac{c\rho_k(N+r)}{N+r}-c\rho_k(N+r)
   =-\frac{r\,c\rho_k(N+r)}{N+r}\longrightarrow0.
\]
Therefore $NS(N)-R(N)\to0$. But the preceding slide gives
$NS(N)=-NE(N)\to0$.
\begin{block}{The consequence of rationality}
$\alpha_k\in\Q$ would force $R(A+Qu)\longrightarrow0$.
\end{block}

\note{[Sources] Local proof: Erdos252/Solution.lean and PROOF.tex. [/Sources]}
\end{frame}

% ----------------------------------------------------------------------

\begin{frame}{One shift has a coefficient that cannot cancel away}

Lemma:

We have
\[
 r_e(j)=r_*:=(k+1)B
\]
\emph{if and only if} $e=0$ and $j=k$.
\begin{itemize}
\item For $j<k$, $r_e(j)\le k(B+G!\,G)=(k+1)B-1$.
\item For $j=k$, $r_e(k)=(k+1)B+
G!\sum_{a<k}(k-a)e_a(k+2)^a$.
All added coefficients are positive, so equality requires $e=0$.
\end{itemize}
Its coefficient is
\[
 c_*=w_0B^{k+1}\stir{k}{k}=B^{k+1}>0,
 \qquad w_0=(-1)^{k(k+1)}=1.
\]

\note{[Sources] Local proof: Erdos252/Solution.lean and PROOF.tex. [/Sources]}
\end{frame}

% ----------------------------------------------------------------------

\begin{frame}{The survivor is concrete when $k=1$}

Recall $p_e=5,7,9$, $t_e=0,2,4$, and $w_e=1,-2,1$.
The progression has modulus $Q=(5\cdot7\cdot9)^2=99225$.
\[
 W(N)=6T_{N/5}-16T_{(N-2)/7}+13T_{(N-4)/9}.
\]
The common first shift ($j=0$) cancels because $5-14+9=0$.
The second shifts have coefficients $25,-98,81$, giving
\[
 \boxed{\ R(N)=25\rho_1(N+10)-98\rho_1(N+12)+81\rho_1(N+14)\ }.
\]
Rationality would force this expression to tend to $0$ as $u\to\infty$.

\note{[Sources] Local proof: Erdos252/Solution.lean and PROOF.tex. [/Sources]}
\end{frame}

% ----------------------------------------------------------------------

\begin{frame}{4. Divisor sums become divisor densities}

Pair complementary divisors to rewrite
\[
 \rho_k(m)=\frac{\sigma_k(m)}{m^k}=\sum_{d\mid m}\frac1{d^k}
 \qquad(m>0).
\]
Along a progression $Qu+a$ with $Q,a>0$, how often does a divisor occur?
\begin{center}
\renewcommand{\arraystretch}{1.3}
\begin{tabular}{@{}c|cccc@{}}
Example $Qu+a=6u+3$ & $d=2$ & $d=3$ & $d=5$ & $d=9$\\
\hline
Density of $d\mid6u+3$ & $0$ & $1$ & $1/5$ & $1/3$
\end{tabular}
\end{center}
In general,
\[
 \text{density of }d\mid Qu+a=
 \begin{cases}\gcd(d,Q)/d,&\gcd(d,Q)\mid a,\\0,&\gcd(d,Q)\nmid a.\end{cases}
\]

\note{[Sources] Local proof: Erdos252/Solution.lean and PROOF.tex. [/Sources]}
\end{frame}

% ----------------------------------------------------------------------

\begin{frame}{4. The actual progression mean is explicit}

For fixed $k\ge1$ and $Q,a>0$,
\[
 \frac1X\sum_{u=0}^{X-1}\rho_k(Qu+a)
 \longrightarrow
 \mu_Q(a):=\sum_{\substack{d\ge1\\\gcd(d,Q)\mid a}}
       \frac{\gcd(d,Q)}{d^{k+1}}\ge1.
\]
Why may we sum the individual divisor densities?
Counted values $Qu+a$ are distinct positive multiples of $d$, so for $X\ge1$,
\[
 \frac1{d^k}\frac{\#\{u<X:d\mid Qu+a\}}X
 \le\frac{QX+a}{X\,d^{k+1}}
 \le\frac{Q+a}{d^{k+1}}.
\]
This bound is summable in $d$. Dominated convergence applies.

\note{[Sources] Local proof: divTerm_cesaro, affine_count_le_quotient,
divTerm_cesaro_bound, tendsto_progMean. [/Sources] Domination is on averaged counts,
not a uniform bound on rho_1, which is unbounded. The term d=1 proves the mean is at least one.}
\end{frame}

% ----------------------------------------------------------------------

\begin{frame}{4. First split the old mean at $L$}

Take a prime $L\nmid Q$. In $Qu+a$, keep the indices $u=Lw+v$,
where $0\le v<L$:
\[
 Q(Lw+v)+a=QLw+a',\qquad a'=Qv+a\equiv a\pmod Q.
\]
In the \emph{old} mean $\mu_Q(a)$, write the divisors containing $L$ as $d=Ld'$.
Since $\gcd(Ld',Q)=\gcd(d',Q)$, their total contribution is
\[
 \sum_{\substack{d'\ge1\\\gcd(d',Q)\mid a}}
       \frac{\gcd(d',Q)}{(Ld')^{k+1}}
 =\frac{\mu_Q(a)}{L^{k+1}}.
\]
Subtracting from the full mean, the divisors with $L\nmid d$ contribute
\[
 \boxed{\ \mu_Q(a)\bigl(1-L^{-(k+1)}\bigr)\ }.
\]

\note{[Sources] Local proof: progMean_multiples and progMean_refine;
PROOF.tex. [/Sources] The index d' ranges over all positive integers, including
multiples of L. Since gcd(d,Q) divides Q, it divides a' exactly when it divides a.}
\end{frame}

% ----------------------------------------------------------------------

\begin{frame}{4. Refinement gives the claimed factor}

Let's compare $\mu_{QL}(a')$ to $\mu_{Q}(a)$.

For divisors with $L\nmid d$, the gcd and compatibility condition stay the same.
For $L\mid d$, the gcd gains one factor of $L$:
\[
 \gcd(d,QL)=L\gcd(d,Q).
\]
\begin{itemize}
\item If $L\nmid a'$, no $QLw+a'$ is divisible by $L$:
these divisor contributions disappear.
\item If $L\mid a'$, each old compatible contribution is multiplied by $L$.
Their total becomes $L\cdot\mu_Q(a)/L^{k+1}=\mu_Q(a)L^{-k}$.
\end{itemize}
Adding this to the unchanged part gives
\[
 \boxed{\ \mu_{QL}(a')=
 \mu_Q(a)\left(1-L^{-(k+1)}+\ind{L\mid a'}L^{-k}\right)\ }.
\]

\note{[Sources] Local proof: meanTerm_refine and progMean_refine;
PROOF.tex. [/Sources] Write g=gcd(d,Q). The new condition gL divides a'
is equivalent to g divides a and L divides a', because gcd(g,L)=1
and a' is congruent to a modulo Q. Each mean is normalized within its own
progression; there is no additional factor 1/L.}
\end{frame}

% ----------------------------------------------------------------------

\begin{frame}{4. Example: one residue gets the extra term}

Take $k=1$, $Q=6$, $a=3$, and $L=5$.
Two subprogressions of $6u+3$ are
\[
 \begin{aligned}
 u=5w:\quad&6u+3=30w+3,   &&5\nmid3,\\
 u=5w+2:\quad&6u+3=30w+15,&&5\mid15.
 \end{aligned}
\]
The refinement formula gives
\[
 \begin{aligned}
 \mu_{30}(3)&=\mu_6(3)\left(1-\frac1{25}\right)
             =\frac{24}{25}\mu_6(3),\\
 \mu_{30}(15)&=\mu_6(3)\left(1-\frac1{25}+\frac15\right)
              =\frac{29}{25}\mu_6(3).
 \end{aligned}
\]
\takeaway{The second mean is larger by exactly $\mu_6(3)/5>0$.}

\note{[Sources] Local proof: progMean_refine; PROOF.tex. [/Sources]
Here v=0 gives a'=3 and v=2 gives a'=15. This numerical example illustrates
the general refinement formula; it is not an additional assumption.}
\end{frame}

% ----------------------------------------------------------------------

\begin{frame}{4. Choose one residue to miss and one to isolate}

Write $i=(e,j)$, $r_i=r_e(j)$ and $c_i=c_{e,j}$, so
$R(N)=\sum_i c_i\rho_k(N+r_i)$.
Choose a prime $L>\max(Q,\max_i r_i)$ and choose two subprogressions of $N=A+Qu$:
\begin{center}
\renewcommand{\arraystretch}{1.6}
\begin{tabular}{@{}p{5.7cm}p{6.2cm}@{}}
\toprule
\textbf{Miss every shift} & \textbf{Hit only the isolated shift}\\
\midrule
$Qv_0+A\equiv0\pmod L$
& $Qv_1+A+r_*\equiv0\pmod L$\\
$Qv_0+A+r_i\equiv r_i\pmod L$
& $Qv_1+A+r_i\equiv r_i-r_*\pmod L$\\
Never zero: $0<r_i<L$.
& Zero exactly when $r_i=r_*$.\\
\bottomrule
\end{tabular}
\end{center}
These give $N=Q(Lw+v_0)+A$ and $N=Q(Lw+v_1)+A$, both in $N=Qu+A$.

\note{[Sources] Local proof: Erdos252/Solution.lean and PROOF.tex. [/Sources]}
\end{frame}

% ----------------------------------------------------------------------

\begin{frame}{4. Compute the first mean, $m_0$}

For $b=0,1$, average along the chosen subprogression $Q(Lw+v_b)+A$:
\[
 \begin{aligned}
 m_b&:=\lim_{X\to\infty}\frac1X\sum_{w=0}^{X-1}R\bigl(Q(Lw+v_b)+A\bigr)\\
    &=\sum_i c_i\mu_{QL}(Qv_b+A+r_i).
 \end{aligned}
\]
Expand the finite sum $R(N)=\sum_i c_i\rho_k(N+r_i)$ and apply the mean
theorem term by term; every offset is positive.

The residue $v_0$ misses every shift: $L\nmid Qv_0+A+r_i$.
So each indicator in the refinement formula is zero:
\[
 \begin{aligned}
 \mu_{QL}(Qv_0+A+r_i)&=\mu_Q(A+r_i)\bigl(1-L^{-(k+1)}\bigr),\\
 m_0&=(1-L^{-(k+1)})\sum_i c_i\mu_Q(A+r_i).
 \end{aligned}
\]

\note{[Sources] Local proof: tendsto_progMean, progMean_refine, and
isolated_shift_not_tendsto_zero; Erdos252/Solution.lean and PROOF.tex. [/Sources]
For each i, Q(Lw+v_b)+A+r_i=QLw+(Qv_b+A+r_i). Thus the new modulus is QL
and the new offset is Qv_b+A+r_i, congruent to A+r_i modulo Q.
Use a=A+r_i and a'=Qv_b+A+r_i in the refinement formula.
Finite linearity allows the coefficients and the sum over i outside the limit,
even though some coefficients are negative. The average uses X terms of the
chosen subprogression, so there is no additional factor 1/L.}
\end{frame}

% ----------------------------------------------------------------------

\begin{frame}{4. Compute $m_1$: exactly one extra term remains}

For $v_1$, $L\mid Qv_1+A+r_i$ holds exactly when $r_i=r_*$.
Substitute these indicators into the same sum of means:
\[
 \begin{aligned}
 m_1&=\sum_i c_i\mu_Q(A+r_i)
       \bigl(1-L^{-(k+1)}+\ind{r_i=r_*}L^{-k}\bigr)\\
    &=m_0+c_*\mu_Q(A+r_*)L^{-k}.
 \end{aligned}
\]
The common part is $m_0$; the indicator selects the single term at $r_*$.
Their difference satisfies
\[
 \alert{m_1-m_0=c_*\mu_Q(A+r_*)L^{-k}>0},
\]
because $c_*>0$, $\mu_Q(A+r_*)\ge1$, and $L^{-k}>0$.
\begin{block}{Contradiction}
Rationality forces $R(A+Qu)\to0$, so each subprogression and its average
tend to zero: $m_0=m_1=0$. This contradicts the positive difference.
Hence $\alpha_k$ is irrational for every $k\ge1$.
\end{block}

\note{[Sources] Local proof: isolated_shift_not_tendsto_zero and
tendsto_survivor_of_rational; Erdos252/Solution.lean and PROOF.tex. [/Sources]
The indicator equals one for exactly one index, since the shift r_* is isolated.
All other terms, including their signs, have the same common factor as in m_0.
The mean of a sequence tending to zero also tends to zero, which gives the
incompatible equalities m_0=m_1=0 under the rationality assumption.}
\end{frame}

% ----------------------------------------------------------------------

\begin{frame}{Degree zero returns to the short tail argument}

For $k=0$, the first denominator term and the same error estimate give
\[
 T_n=\frac{\sigma_0(n+1)}{n+1}+E_n.
\]
The divisor count is positive and at most $64\sqrt{n+1}$, so
\[
 0<T_n\le\frac{64}{\sqrt{n+1}}+O(n^{-3/2})\longrightarrow0.
\]
Rationality would make these tails integers for all sufficiently large $n$.
That is impossible once $0<T_n<1$.
\takeaway{The separate zero-degree argument completes every natural exponent.}

\note{[Sources] Local proof: Erdos252/Solution.lean and PROOF.tex. [/Sources]}
\end{frame}

% ----------------------------------------------------------------------

\begin{frame}{Appendix: the finite Stirling bound used in Lean}

\small
Put
\[
 P_{j,L}(x)=\sum_{i=0}^{L-1}\stir{i}{j}x^{-(i+1)},\qquad
 F_{j,L}(x)=\frac1{x(x-1)\cdots(x-j)}-P_{j,L}(x).
\]
For $x>j+1$, the exact recurrence is
\[
 F_{j+1,L+1}(x)=\frac{F_{j,L}(x)+(j+1)F_{j+1,L}(x)}x.
\]
For $j+1\le H\le x$, induction on $L$ gives
\[
 0\le F_{j,L}(x)\le H^L/x^{L+1}.
\]
At $L=0$, $0<F_{j,0}(x)\le1/x$. In the induction step, the weighted
sum of two errors bounded by $u$ is at most $(j+2)u\le Hu$.
The $j=0$ error is zero after the first order.
\takeaway{This proves the truncated expansion directly, with an explicit error.}

\note{[Sources] Local proof: Erdos252/Solution.lean and PROOF.tex. [/Sources]}
\end{frame}

% ----------------------------------------------------------------------

\begin{frame}{Appendix: scale each remainder before taking limits}

\small
For $x\ge n+1$, the divisor bound gives
\[
 \sigma_k(x)\le64x^k\sqrt x
 \le\frac{64x^{k+1}}{\sqrt{n+1}}.
\]
For a retained term, take $x=n+j+1$ and $n\ge k+1$. Scale first:
\[
 \begin{aligned}
 0\le(n+1)\sigma_k(x)F_{j,k+1}(x)
 &\le x\frac{64x^{k+1}}{\sqrt{n+1}}\frac{(k+1)^{k+1}}{x^{k+2}}\\
 &=\frac{64(k+1)^{k+1}}{\sqrt{n+1}}\longrightarrow0.
 \end{aligned}
\]
There are only $k+1$ retained terms, so their scaled errors sum to a null sequence.
For the omitted tail $O_n$, factorial denominators give
\[
 0\le(n+1)O_n\le\frac{128(k+1)^{k+1}}{\sqrt{n+1}}\longrightarrow0.
\]
Adding the two limits proves $(n+1)E_n\to0$ directly.

\note{[Sources] Local proof: Erdos252/Solution.lean and PROOF.tex. [/Sources]}
\end{frame}

% ----------------------------------------------------------------------

\begin{frame}{Appendix: a compact notation guide}

\small
\begin{center}
\renewcommand{\arraystretch}{1.35}
\begin{tabular}{@{}ll@{}}
\toprule
$k\ge1$ & fixed divisor-power exponent in the grid argument\\
$e\in\{0,\ldots,k+1\}^k$ & one vertex of the finite grid\\
$p_e,\ t_e$ & multiplier and offset for that vertex\\
$r_e(j)=(j+1)p_e-t_e$ & shift of the $(j+1)$-st tail term\\
$w_e=\prod_{a<k}(-1)^{k+1-e_a}\binom{k+1}{e_a}$ & integer finite-difference weight\\
$A,\ Q=\prod_e p_e^2$ & CRT starting point and progression modulus\\
$n_e=(N-t_e)/p_e$ & tail index, for large $N=A+Qu$\\
$c_{e,j}=w_ep_e^{k+1}\stir{k}{j}$ & coefficient after cancellation\\
$r_*=(k+1)B,\ c_*=B^{k+1}$ & the isolated shift and its coefficient\\
$L,\ v_0,\ v_1$ & fresh prime and the two separating residues\\
\bottomrule
\end{tabular}
\end{center}

\note{[Sources] Local proof: Erdos252/Solution.lean and PROOF.tex. [/Sources]}
\end{frame}

% ----------------------------------------------------------------------

\begin{frame}{Appendix: four quantities with different roles}

\small
\begin{align*}
 T_n&=n!\sum_{m>n}\frac{\sigma_k(m)}{m!}
 &&\text{one scaled tail},\\[3pt]
 W(N)&=\sum_e w_e\sigma_k(p_e)T_{n_e}
 &&\text{the weighted integer},\\[3pt]
 S(N)&=\sum_e\sum_{j=0}^k
     \frac{c_{e,j}\rho_k(N+r_e(j))}{N+r_e(j)}
 &&\text{the surviving divided sum},\\[3pt]
 R(N)&=\sum_e\sum_{j=0}^k c_{e,j}\rho_k(N+r_e(j))
 &&\text{the raw survivor}.
\end{align*}
For sufficiently large $N=A+Qu$, rationality gives the chain
\[
 W=S+o(1/N),\quad W\in\Z,\quad W\to0
 \ \Longrightarrow\ W=0\text{ eventually}
 \ \Longrightarrow\ NS\to0
 \ \Longrightarrow\ R\to0.
\]
The last implication uses $NS-R\to0$, not the false identity $NS=R$.

\note{[Sources] Local proof: Erdos252/Solution.lean and PROOF.tex. [/Sources]}
\end{frame}

% ----------------------------------------------------------------------

\begin{frame}{Sources and the complete proof}

\small
\begin{itemize}
\item \textbf{Formal proof and full mathematical exposition}\\
Tokengrinder, \emph{Erd\H{o}s 252} (2026).
\href{https://github.com/tokengr1nder/Erdos252}{github.com/tokengr1nder/Erdos252}\\
\lean{PROOF.pdf} contains all 89 named proof results and their proofs,
38 definitions, and the six statement-audit results.
\item \textbf{Historical context}\\
K. Pratt, \emph{The irrationality of a divisor function series of
Erd\H{o}s and Kac} (2022), introduction.\\
\href{https://arxiv.org/abs/2209.11124}{arXiv:2209.11124}.
\item J. B. Friedlander, F. Luca, M. Stoiciu, \emph{On the irrationality of
a divisor function series}, Integers 7 (2007), A31.\\
\href{https://math.colgate.edu/~integers/h31/h31.pdf}{Original paper}.
\item \textbf{Problem catalogue}\\
T. F. Bloom, \href{https://www.erdosproblems.com/252}{Erd\H{o}s Problem 252}.
\end{itemize}

\note{[Sources] https://arxiv.org/pdf/2209.11124;
https://math.colgate.edu/~integers/h31/h31.pdf; local PROOF.tex and Solution.lean.
[/Sources] Source provenance reviewed 13 September 2026. The direct catalogue page
returned HTTP 403 during this review; its indexed catalogue entry was visible.
No claim of catalogue acceptance is made.}
\end{frame}
\end{document}
